\section{The Self-Modification Taxonomy}

\label{sec:taxonomy}

We organize agent self-modification into four tiers based on scope, persistence,
and trust level. Each tier has characteristic modification objects, lifecycle
policies, and validation requirements.

\begin{definition}[Modification Tier]
A \emph{modification tier} $T_k$ for $k \in \{1, 2, 3, 4\}$ is a tuple
$T_k = (O_k, L_k, V_k, D_k)$ where $O_k$ is the set of modification objects,
$L_k$ is the lifecycle policy (TTL, test-gated, review-gated), $V_k$ is the
validation predicate, and $D_k$ is the deployment mechanism.
\end{definition}

Table~\ref{tab:tiers} summarizes the four tiers.

\begin{table}[H]
\centering
\begin{tabularx}{\textwidth}{l X X X X}
\toprule
Tier & Object & Lifecycle & Validation & Deployment \\
\midrule
1 & Runtime tools & 7-day TTL & Trigger cases & Hot-reload \\
2 & Extensions & Persistent & \texttt{pnpm test} & npm package \\
3 & Skills & Persistent & Peer review & SKILL.md merge \\
4 & Components & Persistent & CI + human PR & Git merge \\
\bottomrule
\end{tabularx}
\caption{The four modification tiers of HHP, ordered by scope and trust level.}
\label{tab:tiers}
\end{table}

\subsection{Tier 1: Runtime Tool Creation}
\label{sec:tier1}

Runtime tool creation is the lightest-weight form of self-modification. When an
agent observes a friction point---a task it cannot perform efficiently with
existing tools---it can generate a new tool on the spot and hot-load it into the
active session.

\begin{definition}[Friction Event]
A \emph{friction event} $\mathcal{F}$ is a tuple $(\tau, e, n)$ where $\tau$ is
the task type, $e$ is the error or inefficiency observed, and $n$ is the number
of times this pattern has recurred in the current session.
\end{definition}

\begin{definition}[Runtime Tool]
A \emph{runtime tool} $\rho$ is a tuple $(\sigma, \pi, \mathcal{C}, t_0, T)$
where $\sigma$ is the tool specification (JSON Schema), $\pi$ is the
implementation (TypeScript function), $\mathcal{C}$ is the set of triggering
failure cases, $t_0$ is the creation timestamp, and $T = 7\text{ days}$ is the
TTL.
\end{definition}

The validation predicate for Tier 1 is:
\begin{equation}
    V_1(\rho) = \bigwedge_{c \in \mathcal{C}} \text{Pass}(\rho.\pi, c),
\end{equation}
where $\text{Pass}(\pi, c)$ is true if and only if running implementation $\pi$
on test case $c$ produces the expected output without error. A generated tool
is only hot-loaded if $V_1(\rho) = \text{true}$.

Hot-loading proceeds via the agent's plugin API:
\begin{verbatim}
// Tool is registered without restarting the session
await agent.tools.register({
  name:        rho.spec.name,
  description: rho.spec.description,
  schema:      rho.spec.inputSchema,
  handler:     rho.implementation,
  ttl:         7 * 24 * 60 * 60 * 1000,   // 7 days in ms
  triggerCases: rho.triggerCases,
});
\end{verbatim}

After the TTL expires, the tool is automatically unregistered. If the tool was
used successfully, a summary is appended to the agent's experience bank for
potential Tier 2 promotion.

\subsection{Tier 2: Extension Development}
\label{sec:tier2}

Extensions are persistent, versioned packages that add durable capability to the
agent harness. An extension consists of a package manifest, an index module
exporting a plugin conforming to the \texttt{BotPlugin} interface, and a test
suite.

\begin{definition}[Extension]
An \emph{extension} $\xi$ is a tuple $(\rho_{\text{pkg}}, \Pi, \Theta, \mathcal{S})$
where $\rho_{\text{pkg}}$ is the package manifest (\texttt{package.json}),
$\Pi$ is the plugin implementation, $\Theta$ is the set of registered tools and
services, and $\mathcal{S}$ is the test suite.
\end{definition}

The plugin API supports four registration types:

\begin{verbatim}
interface BotPlugin {
  name:    string;
  version: string;
  setup(api: PluginAPI): Promise<void>;
}

interface PluginAPI {
  registerTool(spec: ToolSpec, handler: ToolHandler): void;
  registerService(name: string, service: Service): void;
  on(event: AgentEvent, handler: EventHandler): void;
  registerSkill(name: string, skillMd: string): void;
}
\end{verbatim}

The validation predicate for Tier 2 is:
\begin{equation}
    V_2(\xi) = \text{BuildSuccess}(\xi) \wedge \left(\frac{|\mathcal{S}_{\text{pass}}|}{|\mathcal{S}|} = 1\right),
\end{equation}
where $\mathcal{S}_{\text{pass}}$ is the set of passing tests. Extensions must
achieve $100\%$ test passage before registration. Because extensions persist
beyond session TTL, they are deployed via the standard package management
workflow: \texttt{pnpm build}, \texttt{pnpm test}, followed by registration in
the agent's extension registry.

\subsection{Tier 3: Skill Development}
\label{sec:tier3}

Skills encode declarative procedural knowledge---patterns and heuristics that
help agents perform better on recurring task types. Unlike extensions (which
provide executable capabilities), skills provide interpretable, transferable
knowledge that any agent can apply via prompting.

\begin{definition}[Skill]
A \emph{skill} $\varsigma$ is a SKILL.md document with YAML frontmatter and
Markdown body:
\begin{equation}
    \varsigma = \left(\texttt{name}, \texttt{trigger}, \texttt{scope},
    \texttt{tags}, \texttt{body}\right),
\end{equation}
where \texttt{trigger} is a natural language description of when the skill
applies, \texttt{scope} specifies the applicable agent roles, \texttt{tags}
enable retrieval, and \texttt{body} contains the procedural knowledge.
\end{definition}

Skills are stored as Markdown files in the agent's skill library and loaded
into the system prompt for applicable tasks. The validation predicate for Tier 3
is lightweight---skills are reviewed by a coordinating agent before merge, not
executed against test cases. This reflects the fact that skills encode knowledge
rather than functionality, and erroneous skills produce suboptimal behavior
rather than system failures.

\subsection{Tier 4: Cross-Component Hacking}
\label{sec:tier4}

The most powerful and most carefully governed tier covers modifications to the
agent's ecosystem components: MCP tool server, Agent SDK, LLM Gateway, Operative
computer-use system, GRPO training pipelines, and model configuration. These
modifications affect all agents sharing the infrastructure and therefore require
the highest trust level.

\begin{definition}[Cross-Component Modification]
A \emph{cross-component modification} $\mu$ is a tuple
$(\mathcal{R}, \Delta, B, \texttt{pr})$ where $\mathcal{R}$ is the target
repository, $\Delta$ is the git diff representing the proposed change, $B$ is
the isolated git worktree in which $\Delta$ was developed and tested, and
$\texttt{pr}$ is the pull request submitted for human review.
\end{definition}

The worktree isolation mechanism is central to Tier 4 safety and is analyzed
formally in Section~\ref{sec:safety}. The validation predicate is:
\begin{equation}
    V_4(\mu) = \text{CI}(\mu) \wedge \text{HumanApproval}(\mu.\texttt{pr}),
\end{equation}
where $\text{CI}(\mu)$ is the outcome of automated continuous integration and
$\text{HumanApproval}$ is the result of human pull request review. Neither
condition alone is sufficient.

% ============================================================================
% 4. SAFETY FRAMEWORK
% ============================================================================
